\vspace{3pt plus 2pt}\noindent\hspace*{1.2em}\kbvlabel{ಉದಾ}\par\vspace{1pt}

\indent ಭಲೇ ! ಈಜುವ ಪಂದ್ಯದಲ್ಲಿ ಗೆದ್ದವರಲ್ಲಿ ನೀನೇ ಮೊದಲಿಗ !\par

\noindent\hspace*{2.2em}ಅಲ್ಲಿ ನೋಡು ! ವಿಮಾನ ನೆಲಕ್ಕೆ ಅಪ್ಪಳಿಸಿತು !\par

\noindent\hspace*{2.2em}ಈ ಸೀರೆ ಎಷ್ಟು ಅಂದವಾಗಿದೆ !\par

\vspace{4pt plus 4pt minus 1pt}\noindent{\kbvsecfont\fontsize{11pt}{14.5pt}\selectfont\color{kbvaccentdark} (೬)\hspace{0.45em}ಉದ್ಧರಣ ಚಿಹ್ನೆ :}\hspace{0.5em}ನೇರ ಸಂಭಾಷಣೆಯಲ್ಲಿ ಬರುವ ಮಾತು ಗಳನ್ನು ಉದ್ಧರಿಸಿ ಬರೆಯುವಾಗ, ಮತ್ತೊಬ್ಬರು ಹೇಳಿದುದನ್ನು ಅವರು ಹೇಳಿದಂತೆಯೇ ಬರೆಯುವಾಗ, ಅನ್ಯಭಾಷಾಪದಗಳನ್ನು ಬರವಣಿಗೆಯಲ್ಲಿ ತೋರಿಸಿದಾಗ, ಹೆಚ್ಚು ಮಹತ್ವವುಳ್ಳ ಹಾಗೂ ಗಮನ ಸೆಳೆಯುವ ಶಬ್ದ ಗಳನ್ನು ಬರವಣಿಗೆಯಲ್ಲಿ ತಂದಾಗ ಈ ಚಿಹ್ನೆಯನ್ನು ಬಳಸಬೇಕು.\par

\vspace{3pt plus 2pt}\noindent\hspace*{1.2em}\kbvlabel{ಉದಾ}\par\vspace{1pt}

\vspace{2.5pt plus 3pt minus 0.5pt}\noindent\hspace*{1.2em}\kbvnum{(1)}\hspace{0.55em}ಗಾಂಧೀಜಿಯವರು ಗುಂಡಿನೇಟಿಗೆ ಬಲಿಯಾದಾಗ “ಅತಿಯಾದ ಒಳ್ಳೆಯತನವಿರುವುದು ಅಪಾಯಕಾರಿ” ಎಂದು ಜಾರ್ಜ್‌ ಬರ್ನಾರ್ಡ್‌ ಷಾ ಅಭಿಪ್ರಾಯಪಟ್ಟರು.\par

\vspace{2.5pt plus 3pt minus 0.5pt}\noindent\hspace*{1.2em}\kbvnum{(2)}\hspace{0.55em}ಶ್ರೀ ಎಸ್‌. ಎಲ್‌. ಭೈರಪ್ಪನವರ “ಪರ್ವ” ಪ್ರತಿಯೊಬ್ಬ ಕನ್ನಡಿಗನೂ ಓದಬೇಕಾದ ಪುಸ್ತಕ.\par

\vspace{2.5pt plus 3pt minus 0.5pt}\noindent\hspace*{1.2em}\kbvnum{(3)}\hspace{0.55em}ಉಷ್ಣತೆಯನ್ನು ಅಳೆಯುವ ಸಾಧನ “ಥರ್ಮಾಮೀಟರ.”\par

\vspace{2.5pt plus 3pt minus 0.5pt}\noindent\hspace*{1.2em}\kbvnum{(4)}\hspace{0.55em}“ತಾಳಿದವನು ಬಾಳಿಯಾನು” ಎಂಬ ನಾಣ್ಣುಡಿಯನ್ನು ಅರ್ಥ ಮಾಡಿಕೊಂಡರೆ ಜೀವನದಲ್ಲಿ ಸುಖವುಂಟು.\par

\vspace{4pt plus 4pt minus 1pt}\noindent{\kbvsecfont\fontsize{11pt}{14.5pt}\selectfont\color{kbvaccentdark} (೭)\hspace{0.45em}ವಿವರಣಾತ್ಮಕ ಚಿಹ್ನೆ:}\hspace{0.5em}ತಲೆಬರಹದ ಮುಂದೆ ಈ ಚಿಹ್ನೆಯನ್ನು ಹಾಕುವುದುಂಟು; ಆ ತಲೆಬರಹಕ್ಕೆ ಸಂಬಂಧಿಸಿದ ವಿವರಣೆ ಮುಂದೆ ಕಂಡೆಂದು ಇರುವುದು ಎಂಬ ಅಭಿಪ್ರಾಯವನ್ನು ಈ ಚಿಹ್ನೆ ಸೂಚಿಸುವುದು...\par

\noindent\hspace*{1.2em}\kbvlabel{ಉದಾ}\hspace{0.6em}—\par

\indent ವ್ಯಾಕರಣ ಶಾಸ್ತ್ರ : ಶಬ್ದಗಳ ಶುದ್ಧರೂಪವನ್ನು, ಶಬ್ದಗಳಿಗಿರುವ ಸಂಬಂಧವನ್ನು, ಭಾಷೆಯ ಬೆಳವಣಿಗೆಯ ಸ್ವರೂಪವನ್ನು ತಿಳಿಸುವ ಶಾಸ್ತ್ರ;\par

\noindent\hspace*{2.2em}ಗ್ರಂಥವೇ ವ್ಯಾಕರಣ. ಈ ಶಾಸ್ತ್ರದ...\par
